\documentclass[runningheads]{llncs}

\usepackage[T1]{fontenc}
\usepackage[utf8]{inputenc}
\usepackage[spanish]{babel}
\usepackage{graphicx}
\usepackage{booktabs}
\usepackage{amsmath}

\begin{document}

\title{Machete de Metodología Scrum y Fundamentos Ágiles}
\subtitle{Guía Rápida para Estudiantes de Ingeniería}

\author{Martín Gabriel Cazabán}
\institute{Facultad de Ingeniería, Universidad Nacional de Cuyo (UNCuyo) \\
Mendoza, Argentina \\
\email{martincazaban@gmail.com}}

\maketitle

\begin{abstract}
Este documento resume los principios fundamentales de las metodologías ágiles, centrándose en el framework Scrum. Se detallan los roles, eventos y artefactos, así como una breve introducción a las certificaciones Lean Six Sigma (White y Green Belt) aplicadas a la mejora de procesos industriales y de software.

\keywords{Scrum \and Agile \and Lean Six Sigma \and White Belt \and Green Belt.}
\end{abstract}

\section{Metodología Ágil}
La metodología ágil es un enfoque filosófico para la gestión de proyectos que prioriza la adaptabilidad y la entrega de valor continuo. Se basa en el \textbf{Manifiesto Ágil}, cuyos pilares son:
\begin{itemize}
    \item Individuos e interacciones sobre procesos y herramientas.
    \item Software (o producto) funcionando sobre documentación extensiva.
    \item Colaboración con el cliente sobre negociación contractual.
    \item Respuesta ante el cambio sobre seguir un plan.
\end{itemize}

\section{Framework Scrum}
Scrum es un marco de trabajo empírico basado en la transparencia, inspección y adaptación. Es ideal para proyectos de ingeniería donde los requisitos pueden evolucionar.

\subsection{Roles en Scrum}
\begin{enumerate}
    \item \textbf{Product Owner:} Maximiza el valor del producto y gestiona el \textit{Product Backlog}.
    \item \textbf{Scrum Master:} Facilita el proceso, elimina impedimentos y asegura que el equipo siga los valores de Scrum.
    \item \textbf{Developers:} El equipo multidisciplinario encargado de crear el incremento en cada Sprint.
\end{enumerate}

\subsection{Eventos y Artefactos}
El ciclo de vida se organiza en \textbf{Sprints}. Los eventos clave incluyen el \textit{Sprint Planning}, el \textit{Daily Scrum} (15 min), el \textit{Sprint Review} y la \textit{Sprint Retrospective}. Los artefactos principales son el \textit{Product Backlog}, el \textit{Sprint Backlog} y el \textit{Incremento}.

\section{Lean Six Sigma: White Belt y Green Belt}
La integración de Lean Six Sigma en entornos ágiles permite optimizar procesos mediante la reducción de desperdicios y variabilidad.

\begin{table}[]
\centering
\caption{Diferencias entre niveles de certificación Six Sigma.}
\label{tab:belts}
\begin{tabular}{|p{3cm}|p{8cm}|}
\hline
\textbf{Nivel (Belt)} & \textbf{Perfil y Competencias} \\ \hline
\textbf{White Belt} & Nivel inicial. Conoce la terminología básica y apoya en la cultura de mejora continua dentro de la organización. \\ \hline
\textbf{Green Belt} & Nivel intermedio. Aplica herramientas estadísticas y la metodología DMAIC para resolver problemas de calidad y liderar proyectos de mejora. \\ \hline
\end{tabular}
\end{table}

\section{Conclusión}
Para un futuro profesional de la ingeniería, dominar tanto la flexibilidad de Scrum como el rigor analítico de un Green Belt proporciona una ventaja competitiva crítica en la industria 4.0.

\end{document}
